\section{SmartNIC / trusted NIC / secure storage data path}

\begin{frame}[t,shrink=6]{SmartNIC / trusted NIC / secure storage data path：方向开场}
\DeckKicker{方向开场}
\SlideMeta{证据等级}{方向综述}{来源状态：公开材料综合}
{\large\bfseries\color{Ink} SmartNIC / trusted NIC / secure storage data path 的核心是先界定保护对象、管理边界和证据强度，再用三篇主讲材料串起技术演进。\par}\vspace{1.2mm}
\begin{columns}[T,totalwidth=\textwidth]
  \begin{column}{0.45\textwidth}
    \fcolorbox{Rule}{Paper}{%
  \begin{minipage}[t]{0.97\linewidth}
    \vspace{1.1mm}
    \hspace{1.4mm}\begin{minipage}{0.92\linewidth}
      \PanelTitle{内容说明}
      \scriptsize \begin{itemize}
  \item 把 NIC-local root、secure offload、resource management 与 confidential storage 分开写。
  \item 固定三篇主讲材料；`nvidia\_bluefield\_operation\_2025` 与 `khalilov2024osmosis` 作为辅助材料，不计入三篇主讲。
  \item 先讲方向痛点和设计轴，再逐篇说明贡献、限制、证据强度和可落地条件。
  \item 对规范、Survey、draft、vendor evidence 保持显式标注，避免把背景材料写成一手系统证明。
\end{itemize}
    \end{minipage}
    \vspace{1mm}
  \end{minipage}}
  \end{column}
  \begin{column}{0.49\textwidth}
    \fcolorbox{Rule}{Panel}{%
  \begin{minipage}[t]{0.97\linewidth}
    \vspace{1.1mm}
    \hspace{1.4mm}\begin{minipage}{0.92\linewidth}
      \PanelTitle{三篇主讲材料的分工}
      \scriptsize {\tiny
\begin{tabularx}{\linewidth}{@{}YYY@{}}
\toprule
\textbf{论文} & \textbf{定位} & \textbf{证据边界} \\
\midrule
S-NIC: SmartNIC function isolation & 开创/基础入口 & E1 / 同行评审改进证据 \\
TNIC: trusted NIC architecture & 代表性改进 & E1 / 同行评审改进证据 \\
Hazel: secure disaggregated storage data path & 当前 SOTA/补充边界 & E3 / 草案/未批准规范 \\
\bottomrule
\end{tabularx}
}
    \end{minipage}
    \vspace{1mm}
  \end{minipage}}
  \end{column}
\end{columns}
\SourceFootnote{来源：论文原文、官方规范或公开材料｜证据等级：见本页 badge}
\end{frame}


\begin{frame}[t,shrink=6]{S-NIC: SmartNIC function isolation：内容摘要}
\DeckKicker{内容摘要}
\SlideMeta{证据等级}{E1 / 同行评审改进证据}{来源状态：本地 PDF 已核验}
{\large\bfseries\color{Ink} S-NIC 提出 virtual smart NIC，用硬件隔离 SmartNIC 上的 network functions\par}\vspace{1.2mm}
\begin{columns}[T,totalwidth=\textwidth]
  \begin{column}{0.45\textwidth}
    \fcolorbox{Rule}{Paper}{%
  \begin{minipage}[t]{0.97\linewidth}
    \vspace{1.1mm}
    \hspace{1.4mm}\begin{minipage}{0.92\linewidth}
      \PanelTitle{内容说明}
      \scriptsize \begin{itemize}
  \item S-NIC 提出 virtual smart NIC，用硬件隔离 SmartNIC 上的 network functions
  \item sNIC is the foundational peer-reviewed systems paper for cloud SmartNIC isolation.
\end{itemize}
    \end{minipage}
    \vspace{1mm}
  \end{minipage}}
  \end{column}
  \begin{column}{0.49\textwidth}
    \fcolorbox{Rule}{Panel}{%
  \begin{minipage}[t]{0.97\linewidth}
    \vspace{1.1mm}
    \hspace{1.4mm}\begin{minipage}{0.92\linewidth}
      \PanelTitle{证据定位}
      \scriptsize {\tiny
\begin{tabularx}{\linewidth}{@{}YY@{}}
\toprule
\textbf{字段} & \textbf{内容} \\
\midrule
论文定位 & 开创/基础入口 \\
材料类型 & 系统论文 \\
证据等级 & E1 / 同行评审改进证据 \\
来源状态 & 本地 PDF 已核验 \\
\bottomrule
\end{tabularx}
}
    \end{minipage}
    \vspace{1mm}
  \end{minipage}}
  \end{column}
\end{columns}
\SourceFootnote{来源：S-NIC: SmartNIC function isolation｜证据等级：E1 / 同行评审改进证据}
\end{frame}


\begin{frame}[t,shrink=6]{S-NIC: SmartNIC function isolation：研究背景}
\DeckKicker{研究背景}
\SlideMeta{证据等级}{E1 / 同行评审改进证据}{来源状态：本地 PDF 已核验}
{\large\bfseries\color{Ink} SmartNIC/DPU 上运行 packet processing、DPI、NAT、storage offload 后，NIC OS 和其他租户。\par}\vspace{1.2mm}
\begin{columns}[T,totalwidth=\textwidth]
  \begin{column}{0.45\textwidth}
    \fcolorbox{Rule}{Paper}{%
  \begin{minipage}[t]{0.97\linewidth}
    \vspace{1.1mm}
    \hspace{1.4mm}\begin{minipage}{0.92\linewidth}
      \PanelTitle{内容说明}
      \scriptsize \begin{itemize}
  \item SmartNIC/DPU 上运行 packet processing、DPI、NAT、storage offload 后，NIC OS 和其他租户函数可能泄漏或篡改 state
\end{itemize}
    \end{minipage}
    \vspace{1mm}
  \end{minipage}}
  \end{column}
  \begin{column}{0.49\textwidth}
    \fcolorbox{Rule}{Panel}{%
  \begin{minipage}[t]{0.97\linewidth}
    \vspace{1.1mm}
    \hspace{1.4mm}\begin{minipage}{0.92\linewidth}
      \PanelTitle{问题边界}
      \scriptsize {\tiny
\begin{tabularx}{\linewidth}{@{}YY@{}}
\toprule
\textbf{维度} & \textbf{说明} \\
\midrule
保护对象 & SmartNIC / trusted NIC / secure storage data path \\
传统瓶颈 & 把 NIC-local root、secure offload、resource management 与 confidential... \\
本文入口 & S-NIC: SmartNIC function isolation \\
证据限制 & E1 / 同行评审改进证据 \\
\bottomrule
\end{tabularx}
}
    \end{minipage}
    \vspace{1mm}
  \end{minipage}}
  \end{column}
\end{columns}
\SourceFootnote{来源：S-NIC: SmartNIC function isolation｜证据等级：E1 / 同行评审改进证据}
\end{frame}


\begin{frame}[t,shrink=6]{S-NIC: SmartNIC function isolation：解决方案}
\DeckKicker{解决方案}
\SlideMeta{证据等级}{E1 / 同行评审改进证据}{来源状态：本地 PDF 已核验}
{\large\bfseries\color{Ink} 用 memory denylist、locked TLB、cache partitioning、accelerator virtualization、DMA isolation、bus arbitration 和。\par}\vspace{1.2mm}
\begin{columns}[T,totalwidth=\textwidth]
  \begin{column}{0.45\textwidth}
    \fcolorbox{Rule}{Paper}{%
  \begin{minipage}[t]{0.97\linewidth}
    \vspace{1.1mm}
    \hspace{1.4mm}\begin{minipage}{0.92\linewidth}
      \PanelTitle{内容说明}
      \scriptsize \begin{itemize}
  \item 用 memory denylist、locked TLB、cache partitioning、accelerator virtualization、DMA isolation...
\end{itemize}
    \end{minipage}
    \vspace{1mm}
  \end{minipage}}
  \end{column}
  \begin{column}{0.49\textwidth}
    \fcolorbox{Rule}{Panel}{%
  \begin{minipage}[t]{0.97\linewidth}
    \vspace{1.1mm}
    \hspace{1.4mm}\begin{minipage}{0.92\linewidth}
      \PanelTitle{核心设计拆解}
      \scriptsize \begin{itemize}
  \item 01. Virtual SmartNIC launch/destroy lifecycle
  \item 02. Memory denylist, locked TLB, and DMA isolation
  \item 03. Cache partitioning, accelerator virtualization, and bus arbitration
\end{itemize}
    \end{minipage}
    \vspace{1mm}
  \end{minipage}}
  \end{column}
\end{columns}
\SourceFootnote{来源：S-NIC: SmartNIC function isolation｜证据等级：E1 / 同行评审改进证据}
\end{frame}


\begin{frame}[t,shrink=6]{S-NIC: SmartNIC function isolation：实验结果}
\DeckKicker{实验结果}
\SlideMeta{证据等级}{E1 / 同行评审改进证据}{来源状态：本地 PDF 已核验}
{\large\bfseries\color{Ink} 论文报告硬件面积约 +8.89\%，功耗约 +11.45\%，function throughput worst-case 下降小于 1.7\%\par}\vspace{1.2mm}
\begin{columns}[T,totalwidth=\textwidth]
  \begin{column}{0.45\textwidth}
    \fcolorbox{Rule}{Paper}{%
  \begin{minipage}[t]{0.97\linewidth}
    \vspace{1.1mm}
    \hspace{1.4mm}\begin{minipage}{0.92\linewidth}
      \PanelTitle{内容说明}
      \scriptsize \begin{itemize}
  \item 论文报告硬件面积约 +8.89\%，功耗约 +11.45\%，function throughput worst-case 下降小于 1.7\%
\end{itemize}
    \end{minipage}
    \vspace{1mm}
  \end{minipage}}
  \end{column}
  \begin{column}{0.49\textwidth}
    \fcolorbox{Rule}{Panel}{%
  \begin{minipage}[t]{0.97\linewidth}
    \vspace{1.1mm}
    \hspace{1.4mm}\begin{minipage}{0.92\linewidth}
      \PanelTitle{实验/证据边界}
      \scriptsize {\tiny
\begin{tabularx}{\linewidth}{@{}YYY@{}}
\toprule
\textbf{对象} & \textbf{可支撑} & \textbf{边界} \\
\midrule
数据/来源 & S-NIC: SmartNIC function isolation & 本地 PDF 已核验 \\
Baseline/对照 & 以论文或规范原文为准 & 不跨材料外推 \\
指标/对象 & 只使用当前来源可证明的信息 & E1 / 同行评审改进证据 \\
使用边界 & 可进入本方向演进图 & 不能替代更强证据 \\
\bottomrule
\end{tabularx}
}
    \end{minipage}
    \vspace{1mm}
  \end{minipage}}
  \end{column}
\end{columns}
\SourceFootnote{来源：S-NIC: SmartNIC function isolation｜证据等级：E1 / 同行评审改进证据}
\end{frame}


\begin{frame}[t,shrink=6]{S-NIC: SmartNIC function isolation：文章评价}
\DeckKicker{文章评价}
\SlideMeta{证据等级}{E1 / 同行评审改进证据}{来源状态：本地 PDF 已核验}
{\large\bfseries\color{Ink} 为 secure vNIC/vSwitch/offload 提供清晰的 NIC-local resource ownership baseline\par}\vspace{1.2mm}
\begin{columns}[T,totalwidth=\textwidth]
  \begin{column}{0.45\textwidth}
    \fcolorbox{Rule}{Paper}{%
  \begin{minipage}[t]{0.97\linewidth}
    \vspace{1.1mm}
    \hspace{1.4mm}\begin{minipage}{0.92\linewidth}
      \PanelTitle{内容说明}
      \scriptsize \begin{itemize}
  \item 设计优势：为 secure vNIC/vSwitch/offload 提供清晰的 NIC-local resource ownership baseline。
  \item 局限性：没有真实 production NIC silicon，也不是完整 VM/Realm trusted I/O 或 SPDM/TDISP lifecycle。
  \item 商业化潜力：可指导 DPU tenant network-function isolation；落地风险在资源利用率、function chaining 和 vendor...
  \item 证据边界：E1 / Peer-reviewed SOTA；本地 PDF 已核验。
\end{itemize}
    \end{minipage}
    \vspace{1mm}
  \end{minipage}}
  \end{column}
  \begin{column}{0.49\textwidth}
    \fcolorbox{Rule}{Panel}{%
  \begin{minipage}[t]{0.97\linewidth}
    \vspace{1.1mm}
    \hspace{1.4mm}\begin{minipage}{0.92\linewidth}
      \PanelTitle{文章评价}
      \scriptsize {\tiny
\begin{tabularx}{\linewidth}{@{}YY@{}}
\toprule
\textbf{维度} & \textbf{判断} \\
\midrule
设计优势 & 为 secure vNIC/vSwitch/offload 提供清晰的 NIC-local resource ownership baseline。 \\
主要局限 & 没有真实 production NIC silicon，也不是完整 VM/Realm trusted I/O 或 SPDM/TDISP lifecycle。 \\
商业落地 & 可指导 DPU tenant network-function isolation；落地风险在资源利用率、function chaining 和。 \\
讲稿定位 & 开创/基础入口; 证据强度 E1 \\
\bottomrule
\end{tabularx}
}
    \end{minipage}
    \vspace{1mm}
  \end{minipage}}
  \end{column}
\end{columns}
\SourceFootnote{来源：S-NIC: SmartNIC function isolation｜证据等级：E1 / 同行评审改进证据}
\end{frame}


\begin{frame}[t,shrink=6]{TNIC: trusted NIC architecture：内容摘要}
\DeckKicker{内容摘要}
\SlideMeta{证据等级}{E1 / 同行评审改进证据}{来源状态：本地 PDF 已核验}
{\large\bfseries\color{Ink} TNIC 在 NIC 级别构建 minimal silicon root-of-trust，为可信分布式系统提供 transferable authentication 和 non-equivocation\par}\vspace{1.2mm}
\begin{columns}[T,totalwidth=\textwidth]
  \begin{column}{0.45\textwidth}
    \fcolorbox{Rule}{Paper}{%
  \begin{minipage}[t]{0.97\linewidth}
    \vspace{1.1mm}
    \hspace{1.4mm}\begin{minipage}{0.92\linewidth}
      \PanelTitle{内容说明}
      \scriptsize \begin{itemize}
  \item TNIC 在 NIC 级别构建 minimal silicon root-of-trust，为可信分布式系统提供 transferable authentication 和。
  \item TNIC is a peer-reviewed ASPLOS 2025 SOTA trusted NIC system.
\end{itemize}
    \end{minipage}
    \vspace{1mm}
  \end{minipage}}
  \end{column}
  \begin{column}{0.49\textwidth}
    \fcolorbox{Rule}{Panel}{%
  \begin{minipage}[t]{0.97\linewidth}
    \vspace{1.1mm}
    \hspace{1.4mm}\begin{minipage}{0.92\linewidth}
      \PanelTitle{证据定位}
      \scriptsize {\tiny
\begin{tabularx}{\linewidth}{@{}YY@{}}
\toprule
\textbf{字段} & \textbf{内容} \\
\midrule
论文定位 & 代表性改进 \\
材料类型 & 系统论文 \\
证据等级 & E1 / 同行评审改进证据 \\
来源状态 & 本地 PDF 已核验 \\
\bottomrule
\end{tabularx}
}
    \end{minipage}
    \vspace{1mm}
  \end{minipage}}
  \end{column}
\end{columns}
\SourceFootnote{来源：TNIC: trusted NIC architecture｜证据等级：E1 / 同行评审改进证据}
\end{frame}


\begin{frame}[t,shrink=6]{TNIC: trusted NIC architecture：研究背景}
\DeckKicker{研究背景}
\SlideMeta{证据等级}{E1 / 同行评审改进证据}{来源状态：本地 PDF 已核验}
{\large\bfseries\color{Ink} CPU TEE 在跨节点网络 I/O 中 TCB 大、性能差、异构性高，无法作为统一 network trust substrate\par}\vspace{1.2mm}
\begin{columns}[T,totalwidth=\textwidth]
  \begin{column}{0.45\textwidth}
    \fcolorbox{Rule}{Paper}{%
  \begin{minipage}[t]{0.97\linewidth}
    \vspace{1.1mm}
    \hspace{1.4mm}\begin{minipage}{0.92\linewidth}
      \PanelTitle{内容说明}
      \scriptsize \begin{itemize}
  \item CPU TEE 在跨节点网络 I/O 中 TCB 大、性能差、异构性高，无法作为统一 network trust substrate
\end{itemize}
    \end{minipage}
    \vspace{1mm}
  \end{minipage}}
  \end{column}
  \begin{column}{0.49\textwidth}
    \fcolorbox{Rule}{Panel}{%
  \begin{minipage}[t]{0.97\linewidth}
    \vspace{1.1mm}
    \hspace{1.4mm}\begin{minipage}{0.92\linewidth}
      \PanelTitle{问题边界}
      \scriptsize {\tiny
\begin{tabularx}{\linewidth}{@{}YY@{}}
\toprule
\textbf{维度} & \textbf{说明} \\
\midrule
保护对象 & SmartNIC / trusted NIC / secure storage data path \\
传统瓶颈 & 把 NIC-local root、secure offload、resource management 与 confidential... \\
本文入口 & TNIC: trusted NIC architecture \\
证据限制 & E1 / 同行评审改进证据 \\
\bottomrule
\end{tabularx}
}
    \end{minipage}
    \vspace{1mm}
  \end{minipage}}
  \end{column}
\end{columns}
\SourceFootnote{来源：TNIC: trusted NIC architecture｜证据等级：E1 / 同行评审改进证据}
\end{frame}


\begin{frame}[t,shrink=6]{TNIC: trusted NIC architecture：解决方案}
\DeckKicker{解决方案}
\SlideMeta{证据等级}{E1 / 同行评审改进证据}{来源状态：本地 PDF 已核验}
{\large\bfseries\color{Ink} 在 NIC hardware 上实现最小安全 primitive、kernel-bypass stack 和可验证协议\par}\vspace{1.2mm}
\begin{columns}[T,totalwidth=\textwidth]
  \begin{column}{0.45\textwidth}
    \fcolorbox{Rule}{Paper}{%
  \begin{minipage}[t]{0.97\linewidth}
    \vspace{1.1mm}
    \hspace{1.4mm}\begin{minipage}{0.92\linewidth}
      \PanelTitle{内容说明}
      \scriptsize \begin{itemize}
  \item 在 NIC hardware 上实现最小安全 primitive、kernel-bypass stack 和可验证协议
  \item 用 Tamarin 验证核心性质
\end{itemize}
    \end{minipage}
    \vspace{1mm}
  \end{minipage}}
  \end{column}
  \begin{column}{0.49\textwidth}
    \fcolorbox{Rule}{Panel}{%
  \begin{minipage}[t]{0.97\linewidth}
    \vspace{1.1mm}
    \hspace{1.4mm}\begin{minipage}{0.92\linewidth}
      \PanelTitle{核心设计拆解}
      \scriptsize \begin{itemize}
  \item 01. NIC-level silicon root-of-trust
  \item 02. Transferable authentication and non-equivocation
  \item 03. Tamarin-verified protocol core and kernel-bypass stack
\end{itemize}
    \end{minipage}
    \vspace{1mm}
  \end{minipage}}
  \end{column}
\end{columns}
\SourceFootnote{来源：TNIC: trusted NIC architecture｜证据等级：E1 / 同行评审改进证据}
\end{frame}


\begin{frame}[t,shrink=6]{TNIC: trusted NIC architecture：实验结果}
\DeckKicker{实验结果}
\SlideMeta{证据等级}{E1 / 同行评审改进证据}{来源状态：本地 PDF 已核验}
{\large\bfseries\color{Ink} 论文报告相比 CPU-centric TEE systems 最多 6x 性能提升，硬件 TCB 约 2,114 LoC，并用 Tamarin 验。\par}\vspace{1.2mm}
\begin{columns}[T,totalwidth=\textwidth]
  \begin{column}{0.45\textwidth}
    \fcolorbox{Rule}{Paper}{%
  \begin{minipage}[t]{0.97\linewidth}
    \vspace{1.1mm}
    \hspace{1.4mm}\begin{minipage}{0.92\linewidth}
      \PanelTitle{内容说明}
      \scriptsize \begin{itemize}
  \item 论文报告相比 CPU-centric TEE systems 最多 6x 性能提升，硬件 TCB 约 2,114 LoC，并用 Tamarin 验证核心性质
\end{itemize}
    \end{minipage}
    \vspace{1mm}
  \end{minipage}}
  \end{column}
  \begin{column}{0.49\textwidth}
    \fcolorbox{Rule}{Panel}{%
  \begin{minipage}[t]{0.97\linewidth}
    \vspace{1.1mm}
    \hspace{1.4mm}\begin{minipage}{0.92\linewidth}
      \PanelTitle{实验/证据边界}
      \scriptsize {\tiny
\begin{tabularx}{\linewidth}{@{}YYY@{}}
\toprule
\textbf{对象} & \textbf{可支撑} & \textbf{边界} \\
\midrule
数据/来源 & TNIC: trusted NIC architecture & 本地 PDF 已核验 \\
Baseline/对照 & 以论文或规范原文为准 & 不跨材料外推 \\
指标/对象 & 只使用当前来源可证明的信息 & E1 / 同行评审改进证据 \\
使用边界 & 可进入本方向演进图 & 不能替代更强证据 \\
\bottomrule
\end{tabularx}
}
    \end{minipage}
    \vspace{1mm}
  \end{minipage}}
  \end{column}
\end{columns}
\SourceFootnote{来源：TNIC: trusted NIC architecture｜证据等级：E1 / 同行评审改进证据}
\end{frame}


\begin{frame}[t,shrink=6]{TNIC: trusted NIC architecture：文章评价}
\DeckKicker{文章评价}
\SlideMeta{证据等级}{E1 / 同行评审改进证据}{来源状态：本地 PDF 已核验}
{\large\bfseries\color{Ink} 为 network endpoint root-of-trust 提供 peer-reviewed SOTA 和形式化验证证据\par}\vspace{1.2mm}
\begin{columns}[T,totalwidth=\textwidth]
  \begin{column}{0.45\textwidth}
    \fcolorbox{Rule}{Paper}{%
  \begin{minipage}[t]{0.97\linewidth}
    \vspace{1.1mm}
    \hspace{1.4mm}\begin{minipage}{0.92\linewidth}
      \PanelTitle{内容说明}
      \scriptsize \begin{itemize}
  \item 设计优势：为 network endpoint root-of-trust 提供 peer-reviewed SOTA 和形式化验证证据。
  \item 局限性：不保护 workload memory，也不替代 SPDM/TDISP、IOMMU、link encryption 或 device lifecycle 标准。
  \item 商业化潜力：可成为 SmartNIC local root 的学术蓝图；风险在 silicon provisioning、证书链和云厂商互操作。
  \item 证据边界：E1 / Peer-reviewed SOTA；本地 PDF 已核验。
\end{itemize}
    \end{minipage}
    \vspace{1mm}
  \end{minipage}}
  \end{column}
  \begin{column}{0.49\textwidth}
    \fcolorbox{Rule}{Panel}{%
  \begin{minipage}[t]{0.97\linewidth}
    \vspace{1.1mm}
    \hspace{1.4mm}\begin{minipage}{0.92\linewidth}
      \PanelTitle{文章评价}
      \scriptsize {\tiny
\begin{tabularx}{\linewidth}{@{}YY@{}}
\toprule
\textbf{维度} & \textbf{判断} \\
\midrule
设计优势 & 为 network endpoint root-of-trust 提供 peer-reviewed SOTA 和形式化验证证据。 \\
主要局限 & 不保护 workload memory，也不替代 SPDM/TDISP、IOMMU、link encryption 或 device lifecycle 标准。 \\
商业落地 & 可成为 SmartNIC local root 的学术蓝图；风险在 silicon provisioning、证书链和云厂商互操作。 \\
讲稿定位 & 代表性改进; 证据强度 E1 \\
\bottomrule
\end{tabularx}
}
    \end{minipage}
    \vspace{1mm}
  \end{minipage}}
  \end{column}
\end{columns}
\SourceFootnote{来源：TNIC: trusted NIC architecture｜证据等级：E1 / 同行评审改进证据}
\end{frame}


\begin{frame}[t,shrink=6]{Hazel: secure disaggregated storage data path：内容摘要}
\DeckKicker{内容摘要}
\SlideMeta{证据等级}{E3 / 草案/未批准规范}{来源状态：本地 PDF 已核验}
{\large\bfseries\color{Ink} Hazel 为 NVMe-oF disaggregated storage 提供 confidentiality、integrity 和 freshness，并利用 BlueField-3 offload 降低开销\par}\vspace{1.2mm}
\begin{columns}[T,totalwidth=\textwidth]
  \begin{column}{0.45\textwidth}
    \fcolorbox{Rule}{Paper}{%
  \begin{minipage}[t]{0.97\linewidth}
    \vspace{1.1mm}
    \hspace{1.4mm}\begin{minipage}{0.92\linewidth}
      \PanelTitle{内容说明}
      \scriptsize \begin{itemize}
  \item Hazel 为 NVMe-oF disaggregated storage 提供 confidentiality、integrity 和 freshness，并利用。
  \item Hazel is retained as a draft SOTA item for trusted NIC and secure storage data paths.
\end{itemize}
    \end{minipage}
    \vspace{1mm}
  \end{minipage}}
  \end{column}
  \begin{column}{0.49\textwidth}
    \fcolorbox{Rule}{Panel}{%
  \begin{minipage}[t]{0.97\linewidth}
    \vspace{1.1mm}
    \hspace{1.4mm}\begin{minipage}{0.92\linewidth}
      \PanelTitle{证据定位}
      \scriptsize {\tiny
\begin{tabularx}{\linewidth}{@{}YY@{}}
\toprule
\textbf{字段} & \textbf{内容} \\
\midrule
论文定位 & 当前 SOTA/补充边界 \\
材料类型 & 系统论文 \\
证据等级 & E3 / 草案/未批准规范 \\
来源状态 & 本地 PDF 已核验 \\
\bottomrule
\end{tabularx}
}
    \end{minipage}
    \vspace{1mm}
  \end{minipage}}
  \end{column}
\end{columns}
\SourceFootnote{来源：Hazel: secure disaggregated storage data path｜证据等级：E3 / 草案/未批准规范}
\end{frame}


\begin{frame}[t,shrink=6]{Hazel: secure disaggregated storage data path：研究背景}
\DeckKicker{研究背景}
\SlideMeta{证据等级}{E3 / 草案/未批准规范}{来源状态：本地 PDF 已核验}
{\large\bfseries\color{Ink} confidential workload 使用远端存储时，本地 dm-crypt/dm-integrity/dm-x 难以保持性能、扩展性和 replay freshness\par}\vspace{1.2mm}
\begin{columns}[T,totalwidth=\textwidth]
  \begin{column}{0.45\textwidth}
    \fcolorbox{Rule}{Paper}{%
  \begin{minipage}[t]{0.97\linewidth}
    \vspace{1.1mm}
    \hspace{1.4mm}\begin{minipage}{0.92\linewidth}
      \PanelTitle{内容说明}
      \scriptsize \begin{itemize}
  \item confidential workload 使用远端存储时，本地 dm-crypt/dm-integrity/dm-x 难以保持性能、扩展性和 replay freshness
\end{itemize}
    \end{minipage}
    \vspace{1mm}
  \end{minipage}}
  \end{column}
  \begin{column}{0.49\textwidth}
    \fcolorbox{Rule}{Panel}{%
  \begin{minipage}[t]{0.97\linewidth}
    \vspace{1.1mm}
    \hspace{1.4mm}\begin{minipage}{0.92\linewidth}
      \PanelTitle{问题边界}
      \scriptsize {\tiny
\begin{tabularx}{\linewidth}{@{}YY@{}}
\toprule
\textbf{维度} & \textbf{说明} \\
\midrule
保护对象 & SmartNIC / trusted NIC / secure storage data path \\
传统瓶颈 & 把 NIC-local root、secure offload、resource management 与 confidential... \\
本文入口 & Hazel: secure disaggregated storage data path \\
证据限制 & E3 / 草案/未批准规范 \\
\bottomrule
\end{tabularx}
}
    \end{minipage}
    \vspace{1mm}
  \end{minipage}}
  \end{column}
\end{columns}
\SourceFootnote{来源：Hazel: secure disaggregated storage data path｜证据等级：E3 / 草案/未批准规范}
\end{frame}


\begin{frame}[t,shrink=6]{Hazel: secure disaggregated storage data path：解决方案}
\DeckKicker{解决方案}
\SlideMeta{证据等级}{E3 / 草案/未批准规范}{来源状态：本地 PDF 已核验}
{\large\bfseries\color{Ink} counter leasing、NVMe metadata、Hazel Merkle Tree、metadata cache、eventual consistency 和 DPU crypto offload\par}\vspace{1.2mm}
\begin{columns}[T,totalwidth=\textwidth]
  \begin{column}{0.45\textwidth}
    \fcolorbox{Rule}{Paper}{%
  \begin{minipage}[t]{0.97\linewidth}
    \vspace{1.1mm}
    \hspace{1.4mm}\begin{minipage}{0.92\linewidth}
      \PanelTitle{内容说明}
      \scriptsize \begin{itemize}
  \item counter leasing、NVMe metadata、Hazel Merkle Tree、metadata cache、eventual consistency 和 DPU...
\end{itemize}
    \end{minipage}
    \vspace{1mm}
  \end{minipage}}
  \end{column}
  \begin{column}{0.49\textwidth}
    \fcolorbox{Rule}{Panel}{%
  \begin{minipage}[t]{0.97\linewidth}
    \vspace{1.1mm}
    \hspace{1.4mm}\begin{minipage}{0.92\linewidth}
      \PanelTitle{核心设计拆解}
      \scriptsize \begin{itemize}
  \item 01. Counter leasing and KBS/TEE-assisted control path
  \item 02. NVMe metadata plus Hazel Merkle Tree freshness
  \item 03. BlueField-3 crypto offload for NVMe-oF data path
\end{itemize}
    \end{minipage}
    \vspace{1mm}
  \end{minipage}}
  \end{column}
\end{columns}
\SourceFootnote{来源：Hazel: secure disaggregated storage data path｜证据等级：E3 / 草案/未批准规范}
\end{frame}


\begin{frame}[t,shrink=6]{Hazel: secure disaggregated storage data path：实验结果}
\DeckKicker{实验结果}
\SlideMeta{证据等级}{E3 / 草案/未批准规范}{来源状态：本地 PDF 已核验}
{\large\bfseries\color{Ink} arXiv preprint 提供 NVMe-oF storage 原型评估，报告常见路径 1\%-2\% overhead，IO500 平均约 6.。\par}\vspace{1.2mm}
\begin{columns}[T,totalwidth=\textwidth]
  \begin{column}{0.45\textwidth}
    \fcolorbox{Rule}{Paper}{%
  \begin{minipage}[t]{0.97\linewidth}
    \vspace{1.1mm}
    \hspace{1.4mm}\begin{minipage}{0.92\linewidth}
      \PanelTitle{内容说明}
      \scriptsize \begin{itemize}
  \item arXiv preprint 提供 NVMe-oF storage 原型评估，报告常见路径 1\%-2\% overhead，IO500 平均约 6.3\%，YCSB p99...
  \item 小随机 freshness I/O 是弱点
\end{itemize}
    \end{minipage}
    \vspace{1mm}
  \end{minipage}}
  \end{column}
  \begin{column}{0.49\textwidth}
    \fcolorbox{Rule}{Panel}{%
  \begin{minipage}[t]{0.97\linewidth}
    \vspace{1.1mm}
    \hspace{1.4mm}\begin{minipage}{0.92\linewidth}
      \PanelTitle{实验/证据边界}
      \scriptsize {\tiny
\begin{tabularx}{\linewidth}{@{}YYY@{}}
\toprule
\textbf{对象} & \textbf{可支撑} & \textbf{边界} \\
\midrule
数据/来源 & Hazel: secure disaggregated storage data path & 本地 PDF 已核验 \\
Baseline/对照 & 以论文或规范原文为准 & 不跨材料外推 \\
指标/对象 & 只使用当前来源可证明的信息 & E3 / 草案/未批准规范 \\
使用边界 & 可进入本方向演进图 & 不能替代更强证据 \\
\bottomrule
\end{tabularx}
}
    \end{minipage}
    \vspace{1mm}
  \end{minipage}}
  \end{column}
\end{columns}
\SourceFootnote{来源：Hazel: secure disaggregated storage data path｜证据等级：E3 / 草案/未批准规范}
\end{frame}


\begin{frame}[t,shrink=6]{Hazel: secure disaggregated storage data path：文章评价}
\DeckKicker{文章评价}
\SlideMeta{证据等级}{E3 / 草案/未批准规范}{来源状态：本地 PDF 已核验}
{\large\bfseries\color{Ink} 把 confidential computing 的 storage data path、freshness 和 DPU offload 放到同一。\par}\vspace{1.2mm}
\begin{columns}[T,totalwidth=\textwidth]
  \begin{column}{0.45\textwidth}
    \fcolorbox{Rule}{Paper}{%
  \begin{minipage}[t]{0.97\linewidth}
    \vspace{1.1mm}
    \hspace{1.4mm}\begin{minipage}{0.92\linewidth}
      \PanelTitle{内容说明}
      \scriptsize \begin{itemize}
  \item 设计优势：把 confidential computing 的 storage data path、freshness 和 DPU offload 放到同一系统原型中。
  \item 局限性：仍是 preprint；没有把 BlueField DPU attestation、SPDM/TDISP/NVMe-oF endpoint identity 完整闭环。
  \item 商业化潜力：接近云上 confidential storage 需求；风险在 crash consistency、KBS/DPU 运维和 production endpoint...
  \item 证据边界：E3 / Draft-not-ratified；本地 PDF 已核验。
\end{itemize}
    \end{minipage}
    \vspace{1mm}
  \end{minipage}}
  \end{column}
  \begin{column}{0.49\textwidth}
    \fcolorbox{Rule}{Panel}{%
  \begin{minipage}[t]{0.97\linewidth}
    \vspace{1.1mm}
    \hspace{1.4mm}\begin{minipage}{0.92\linewidth}
      \PanelTitle{文章评价}
      \scriptsize {\tiny
\begin{tabularx}{\linewidth}{@{}YY@{}}
\toprule
\textbf{维度} & \textbf{判断} \\
\midrule
设计优势 & 把 confidential computing 的 storage data path、freshness 和 DPU offload 放到同一系统原型中。 \\
主要局限 & 仍是 preprint；没有把 BlueField DPU attestation、SPDM/TDISP/NVMe-oF endpoint... \\
商业落地 & 接近云上 confidential storage 需求；风险在 crash consistency、KBS/DPU 运维和 production... \\
讲稿定位 & 当前 SOTA/补充边界; 证据强度 E3 \\
\bottomrule
\end{tabularx}
}
    \end{minipage}
    \vspace{1mm}
  \end{minipage}}
  \end{column}
\end{columns}
\SourceFootnote{来源：Hazel: secure disaggregated storage data path｜证据等级：E3 / 草案/未批准规范}
\end{frame}


\begin{frame}[t,shrink=6]{SmartNIC / trusted NIC / secure storage data path：方向总结}
\DeckKicker{方向总结}
\SlideMeta{证据等级}{方向综述}{来源状态：公开材料综合}
{\large\bfseries\color{Ink} SmartNIC / trusted NIC / secure storage data path 的结论是：三篇主讲材料共同给出技术演进，但仍要用证据等级约束每个结论。\par}\vspace{1.2mm}
\begin{columns}[T,totalwidth=\textwidth]
  \begin{column}{0.45\textwidth}
    \fcolorbox{Rule}{Paper}{%
  \begin{minipage}[t]{0.97\linewidth}
    \vspace{1.1mm}
    \hspace{1.4mm}\begin{minipage}{0.92\linewidth}
      \PanelTitle{内容说明}
      \scriptsize \begin{itemize}
  \item S-NIC: SmartNIC function isolation：开创/基础入口，E1 / 同行评审改进证据。
  \item TNIC: trusted NIC architecture：代表性改进，E1 / 同行评审改进证据。
  \item Hazel: secure disaggregated storage data path：当前 SOTA/补充边界，E3 / 草案/未批准规范。
  \item 本方向结论：先用三篇主讲材料建立演进关系，再用证据等级控制机制、实验和商业化结论。
\end{itemize}
    \end{minipage}
    \vspace{1mm}
  \end{minipage}}
  \end{column}
  \begin{column}{0.49\textwidth}
    \fcolorbox{Rule}{Panel}{%
  \begin{minipage}[t]{0.97\linewidth}
    \vspace{1.1mm}
    \hspace{1.4mm}\begin{minipage}{0.92\linewidth}
      \PanelTitle{本方向方法对比}
      \scriptsize {\tiny
\begin{tabularx}{\linewidth}{@{}YYY@{}}
\toprule
\textbf{论文} & \textbf{定位} & \textbf{选择理由} \\
\midrule
S-NIC: SmartNIC function isolation & 开创/基础入口 & sNIC is the foundational peer-reviewed systems paper for cloud SmartNIC... \\
TNIC: trusted NIC architecture & 代表性改进 & TNIC is a peer-reviewed ASPLOS 2025 SOTA trusted NIC system. \\
Hazel: secure disaggregated storage data path & 当前 SOTA/补充边界 & Hazel is retained as a draft SOTA item for trusted NIC and secure... \\
\bottomrule
\end{tabularx}
}
    \end{minipage}
    \vspace{1mm}
  \end{minipage}}
  \end{column}
\end{columns}
\SourceFootnote{来源：论文原文、官方规范或公开材料｜证据等级：见本页 badge}
\end{frame}
